\section*{Глава 3. Разработка веб-приложения}
\addcontentsline{toc}{section}{Глава 3. Разработка веб-приложения}
\stepcounter{section}

\subsection{Разработка базовой структуры приложения и вёрстка интерфейса}

Разработка веб-приложения была выполнена на основе фреймворка Django. Для максимальной простоты и удобства понимания весь основной функционал магазина был вынесен в одно приложение \texttt{store}. Это позволяет избежать излишней сложности и легко отслеживать связи между товарами, корзиной и заказами.

Интерфейс спроектирован с использованием базовых HTML-шаблонов и CSS (стили встроены напрямую в базовый шаблон для простоты чтения кода). Пользователю доступны главная страница с каталогом товаров, страница с детальной информацией о товаре, корзина, оформление заказа и история покупок.

\subsection{Проектирование моделей данных (товары и заказы)}

Для хранения информации в базе данных (используется стандартная база SQLite) были созданы понятные модели. Основными являются \texttt{Category} (Категория) и \texttt{Product} (Товар). 

\begin{lstlisting}[caption={Модели товаров и категорий}, language=Python]
class Category(models.Model):
    name = models.CharField(max_length=255, verbose_name="Название")
    description = models.TextField(blank=True)

class Product(models.Model):
    category = models.ForeignKey(Category, on_delete=models.CASCADE)
    name = models.CharField(max_length=255)
    description = models.TextField()
    price = models.DecimalField(max_digits=10, decimal_places=2)
    stock = models.PositiveIntegerField(default=0)
\end{lstlisting}

Связь между ними организована по принципу "один ко многим" (одна категория может содержать много товаров).

\subsection{Реализация корзины товаров}

Корзина товаров реализована наиболее простым и логичным способом — через встроенные сессии Django (\texttt{request.session}). Это значит, что корзина хранится во временной памяти сервера для каждого пользователя (даже неавторизованного) и не требует создания сложных таблиц в базе данных. В сессии хранится словарь, где ключ — это ID товара, а значение — его количество.

\begin{lstlisting}[caption={Добавление товара в корзину}, language=Python]
def add_to_cart(request, product_id):
    if 'cart' not in request.session:
        request.session['cart'] = {}
    cart = request.session['cart']
    
    product_id_str = str(product_id)
    if product_id_str in cart:
        cart[product_id_str] += 1 # увеличиваем количество
    else:
        cart[product_id_str] = 1  # добавляем новый товар
        
    request.session.modified = True
    return redirect('cart_detail')
\end{lstlisting}

\subsection{Оформление заказа и аутентификация}

Для оформления заказа пользователь должен быть авторизован. Мы использовали стандартную модель \texttt{User} из Django, что избавило от необходимости писать систему регистрации с нуля. Авторизованный покупатель может заполнить адрес, после чего содержимое его корзины из сессии переносится в базу данных в виде \texttt{Order} (Заказ) и \texttt{OrderItem} (Позиции заказа). Цены фиксируются на момент покупки.

\begin{lstlisting}[caption={Оформление заказа}, language=Python]
@login_required
def checkout(request):
    cart = request.session.get('cart', {})
    if request.method == 'POST':
        address = request.POST.get('address')
        # Создаем заказ в базе данных
        order = Order.objects.create(user=request.user, address=address)
        # Логика переноса товаров из корзины в позиции заказа...
        request.session['cart'] = {} # Очищаем корзину после заказа
        return redirect('order_history')
    return render(request, 'store/checkout.html')
\end{lstlisting}

\subsection{Реализация поиска и фильтрации товаров}

Для удобства пользователей на главной странице добавлен поиск и фильтрация. Поиск осуществляется по названию товара с использованием простого фильтра \texttt{icontains} (поиск подстроки без учета регистра), а фильтрация — по точному совпадению ID категории.

\begin{lstlisting}[caption={Поиск и фильтрация в каталоге}, language=Python]
def product_list(request):
    query = request.GET.get('q', '')
    category_id = request.GET.get('category', '')
    
    products = Product.objects.all()
    if query: # Если есть поисковый запрос
        products = products.filter(name__icontains=query)
    if category_id: # Если выбрана категория
        products = products.filter(category_id=category_id)
        
    return render(request, 'store/product_list.html', {'products': products})
\end{lstlisting}

Таким образом, веб-приложение реализует весь необходимый функционал интернет-магазина, оставаясь при этом логичным, легко читаемым и удобным для дальнейшего расширения или изучения на любительском уровне.

Исходный код разработанного веб-приложения доступен в репозитории на платформе GitHub по ссылке: \url{https://github.com/DanilaKuznetsov/coursework-web-app-main}. Доступ к репозиторию предоставлен проверяющему (GitHub: \texttt{akruzhalov}).
